\documentclass[11pt]{article}
\usepackage{fontspec}
\setmainfont{DejaVu Serif}
\usepackage{amsmath,amssymb,amsthm,mathtools,mathrsfs,esint}
\usepackage[margin=1in]{geometry}
\usepackage[hidelinks]{hyperref}
\providecommand{\tightlist}{\setlength{\itemsep}{0pt}\setlength{\parskip}{0pt}}
\setlength{\emergencystretch}{3em}
\title{The Weak Endpoint at Every Fixed Finite Euler Density}
\author{Qiushi Line 1 --- continuation C}
\date{29 September 2026}
\begin{document}
\maketitle
\noindent The new conclusion is quantified as: for every finite $B$, there exists $c(B)>0$. It does not assert a density-independent constant and does not determine the optimal coefficient in O11. The main proof is followed by its retained strong-current input, a shorter bounded resolvent construction, the amplitude/phase calculation, and the more general line-bundle pairing lemma. Actual field identities and the finite-parameter/conformal comparison of the supplied manuscript remain explicitly inherited. No Nijenhuis-bracket prolongation is needed in this proof.
\tableofcontents
\clearpage
\hypertarget{o10-closes-at-every-fixed-finite-euler-density}{%
\section{O10 closes at every fixed finite Euler
density}\label{o10-closes-at-every-fixed-finite-euler-density}}

Date: 2026-09-29. This continuation uses the already established
strong-current result, a direct mass--phase--Dirac reduction, the new
constant-length derivative lemma, and a logarithmic resolvent argument
for a line bundle. The conclusion removes the previous numerical
Euler-density threshold. It does not assert a constant independent of
the Euler-density bound.

\hypertarget{precise-conclusion}{%
\subsection{1. Precise conclusion}\label{precise-conclusion}}

Retain the original O10 actual conformal monopole input, the positive
projection budget \(p\ge A>0\), and \(a^2\ge T_+\). For every fixed
finite \(B\), there exists \(c(B)>0\) such that

\[
 \boxed{\frac{\chi(M)}A\le B
 \quad\Longrightarrow\quad
 \fint_M|\operatorname{Ric}+g|^2+\frac{D_\beta(g)}{V(g)}
 \ge c(B)\frac{(\beta-2)^2}{\beta^2},\qquad\beta\ge2.}
\tag{1}
\]

The class of admissible metrics and the actual supplier are exactly
those of the original O10 problem. There is no added pointwise Ricci
bound, diameter bound, noncollapse hypothesis or Sobolev inequality. The
dependence on the fixed number \(B\) is retained.

Here \(a=c_1(L)\), \(p=(a_g^+)^2\),
\(T_+=2\chi(M)+3\operatorname{sign}(M)\), and

\[
 D_\beta=(1+\beta^{-1})\int s_g^2
 +6(\beta+1)\int|W_g^+|^2-72\pi^2p.
\]

Weyl norms are Hilbert--Schmidt norms on two-forms; Ricci norms are
tensor norms. The actual single solution is taken on the negative Yamabe
representative normalized by \(s=-4\). An application to every original
metric uses the original conformal-class supplier and the stated
conformal comparison. Throughout a sequence, all analytic norms are
taken with respect to \(dV_j/V_j\).

Equivalently, an unrestricted hypothetical O10 zero-end sequence must
have

\[
 \boxed{\chi(M_j)/A_j\longrightarrow+\infty}
\tag{2}
\]

along the endpoint subsequence. The statement does not yet rule out this
final escape regime.

\hypertarget{the-new-analytic-lemma-a-line-cannot-keep-characteristic-square-through-this-bounded-energy-limit}{%
\subsection{2. The new analytic lemma: a line cannot keep characteristic
square through this bounded-energy
limit}\label{the-new-analytic-lemma-a-line-cannot-keep-characteristic-square-through-this-bounded-energy-limit}}

Consider closed oriented Riemannian four-manifolds \((M_j,g_j)\),
Hermitian line bundles \(K_j\) and unitary connections. Write the real
curvature as \(F_j=-iF^{K_j}\). Suppose:

\begin{enumerate}
\def\labelenumi{\arabic{enumi}.}
\tightlist
\item
  \(\fint|\operatorname{Ric}_{g_j}|^2\le C_R\).
\item
  \(F_j=F_{0,j}+e_j\), with \(|F_{0,j}|\le C_F\) and
  \(\|e_j\|_{2,\mathrm{avg}}\to0\). The bounded model need not be
  closed.
\item
  Smooth sections \(\sigma_j\) satisfy \(\fint|\nabla\sigma_j|^2\le E\),
  and \(|\sigma_j|\to\ell>0\) strongly in normalized \(L^2\).
\end{enumerate}

Then

\[
 \boxed{\frac{c_1(K_j)^2[M_j]}{V_j}
 =\frac1{4\pi^2}\fint F_j\wedge F_j\longrightarrow0.}
\tag{3}
\]

The notation on the right means the coefficient of the four-form
relative to the volume form. The conclusion is signed convergence; no
sign of the characteristic number is assumed.

\hypertarget{a-bounded-cap-and-its-spectral-regularization}{%
\subsubsection{2.1 A bounded cap and its spectral
regularization}\label{a-bounded-cap-and-its-spectral-regularization}}

Apply a fixed smooth radial cap which is the identity near radius
\(\ell\), has uniformly bounded derivative, and has bounded image. The
resulting sections \(s_j\) satisfy

\[
 |s_j|\le M,\qquad \fint|\nabla s_j|^2\le E',
 \qquad d_j:=\fint(|s_j|-\ell)^2\to0.
\tag{4}
\]

Let \(L_j=\nabla^*\nabla\) and, for \(\varepsilon>0\), put

\[
 u_{j,\varepsilon}=(I+\varepsilon L_j)^{-1}s_j.
\]

The connection-Laplacian maximum principle and the spectral theorem give

\[
 |u_{j,\varepsilon}|\le M,\qquad
 \fint|u_{j,\varepsilon}-s_j|^2\le\varepsilon E',
 \qquad \fint|\nabla u_{j,\varepsilon}|^2\le E'.
\tag{5}
\]

For \(K_j(\varepsilon)=\fint|L_j u_{j,\varepsilon}|^2\), the exact
spectral calculation is

\[
 \int_0^1K_j(\varepsilon)\,d\varepsilon
 =\frac1{V_j}\sum_\nu
 \frac{\lambda_\nu^2}{1+\lambda_\nu}|(s_j)_\nu|^2
 \le\fint|\nabla s_j|^2\le E'.
\tag{6}
\]

No uniform spectral gap is required.

Choose numbers \(\delta_j\downarrow0\) with

\[
 d_j+\|e_j\|_{2,\mathrm{avg}}^2\le\delta_j<1.
\]

Since

\[
 \int_{\delta_j^{1/2}}^{\delta_j^{1/4}}
 K_j(\varepsilon)\,d\varepsilon
 =\int_{\delta_j^{1/2}}^{\delta_j^{1/4}}
 [\varepsilon K_j(\varepsilon)]\frac{d\varepsilon}{\varepsilon}
 \le E',
\]

there exists \(\varepsilon_j\) in this interval such that

\[
 \boxed{\varepsilon_jK_j(\varepsilon_j)
 \le\frac{4E'}{|\log\delta_j|}\longrightarrow0.}
\tag{7}
\]

Write \(u_j=u_{j,\varepsilon_j}\), \(K_j=K_j(\varepsilon_j)\). In
particular \(\varepsilon_j\to0\), while (7) is stronger than the
elementary bound \(K_j\le E'/\varepsilon_j\).

\hypertarget{the-geometry-free-fourth-power-estimate}{%
\subsubsection{2.2 The geometry-free fourth-power
estimate}\label{the-geometry-free-fourth-power-estimate}}

For a bounded section \(u\) of a line bundle over a closed
four-manifold, derivative interchange and integration by parts give

\[
 \fint\bigl(|\nabla^2u|^2+|\nabla u|^4\bigr)
 \le C_M\left\{\fint|Lu|^2+
 \fint\bigl(|\operatorname{Ric}|^2+|F|^2\bigr)\right\}.
\tag{8}
\]

A harmless additive constant can be included. Only the Ricci contraction
occurs when commuting derivatives on the covector slot. Bundle curvature
enters without derivatives. To see the absorption, set
\(I=\int|\nabla^2u|^2\), \(J=\int|\nabla u|^4\) and \(K=\int|Lu|^2\).
Interchanging the two covariant derivatives first gives

\[
 I\le C K+C M^2\int|F|^2+
 C\int(|\operatorname{Ric}|+|F|)|\nabla u|^2.
\]

Integration of
\(\operatorname{div}(|\nabla u|^2\operatorname{Re}\langle u,\nabla u\rangle)\)
gives \(J\le C_M(I+K)\). Cauchy--Schwarz and Young then absorb \(I\) and
yield (8). No Sobolev estimate or derivative of curvature is used.

Thus in our sequence

\[
 \fint|\nabla u_j|^4\le C(1+K_j).
\tag{9}
\]

Let \(Z_j=\{|u_j|\le3\ell/4\}\). From (4)--(5),

\[
 \frac{\operatorname{Vol}(Z_j)}{V_j}
 \le C(d_j+\varepsilon_j E')\le C\varepsilon_j.
\]

Consequently

\[
 \boxed{\fint_{Z_j}|\nabla u_j|^2
 \le C\sqrt{\varepsilon_j(1+K_j)}\longrightarrow0.}
\tag{10}
\]

This is the estimate unavailable for the unregularized \(H^1\) section.

\hypertarget{rank-one-transgression-and-its-error-bounds}{%
\subsubsection{2.3 Rank-one transgression and its error
bounds}\label{rank-one-transgression-and-its-error-bounds}}

Put \(a_j=|u_j|^2\),
\(j_j=\operatorname{Im}\langle u_j,\nabla u_j\rangle\). Choose a fixed
smooth cutoff \(\chi(a)\), zero for \(a\le(\ell/2)^2\) and one for
\(a\ge(3\ell/4)^2\). The real one-form

\[
 \theta_j=\chi(a_j)j_j/a_j
\]

is globally smooth, extended by zero at zeros of \(u_j\). For a line
bundle the exact local identity is

\[
 d\theta_j=\chi(a_j)F_j+
 \chi'(a_j)\,da_j\wedge j_j/a_j.
\tag{11}
\]

Its second term has norm at most \(C_\ell|\nabla u_j|^2\) and is
supported in \(Z_j\). Since \(dF_j=0\), Stokes gives

\[
 \left|\fint F_j\wedge F_j\right|
 \le C\fint_{Z_j}|F_j|^2+
 C\fint_{Z_j}|F_j|\,|\nabla u_j|^2.
\tag{12}
\]

The first term tends to zero because \(F_j\) has a uniformly bounded
model and an \(L^2\)-small error. For the second term, split
\(F_j=F_{0,j}+e_j\). Its bounded-model part tends to zero by (10). Its
error part satisfies

\[
 \fint|e_j|\,|\nabla u_j|^2
 \le C\sqrt{\delta_j(1+K_j)}\longrightarrow0,
\tag{13}
\]

since \(\delta_j/\varepsilon_j\le\delta_j^{1/2}\) and
\(\varepsilon_jK_j\to0\). This proves (3).

\hypertarget{applying-the-lemma-to-the-actual-o10-system}{%
\subsection{3. Applying the lemma to the actual O10
system}\label{applying-the-lemma-to-the-actual-o10-system}}

Suppose that a scalar-minus-four actual endpoint sequence has

\[
 \beta_j\to\infty,\qquad
 \kappa_j=P_j/V_j+D_{\beta_j}/V_j\to0,
 \qquad \chi(M_j)/A_j\le B<\infty.
\tag{14}
\]

The existing curvature/topological budgets imply
\(\fint|W^-|^2\le C_B\). The preceding strong-current and resolvent
arguments produce \(v\in S^-\otimes L\) with

\[
 v\text{ bounded},\quad |v|^2\to m_*,\quad
 \nabla^*\nabla v-2v\to0\text{ in }L^2,
 \quad \|\nabla v\|_4+\|\nabla^2v\|_2\le C_B.
\]

The direct mass--phase--Dirac reduction leaves only
\(C=(\nabla_{\bar Z_2}v)_2\). Indeed, the inherited
\(N_v=(\nabla v)^{1,0}_{J_+}\to0\) gives \(A_i=\nabla_{Z_i}v\to0\). The
scalar identity and bounded second-derivative budget give \(da\to0\),
where \(a=|v|^2\). The exact phase identity
\(|j_v+J_+da/2|^2\le2a|N_v|^2\) gives \(j_v\to0\). In a frame with
\(v=\sqrt a\,e_1\), \[
 (\nabla v)_1=(da/2+i j_v)/\sqrt a,
\] so every first spinor component tends to zero on fixed
positive-amplitude sets. Low-amplitude sets contribute \(o(1)\) to
\(\fint|\nabla v|^2\) by the \(L^4\) bound. Thus \(B_1,C_1\to0\), and
the Clifford identity \(D^-v=((A_1)_1-(A_2)_2,B_2+C_1)\) gives
\(B_2\to0\) as well. No bracket prolongation estimate is required.

The new covariance argument in \texttt{O10\_\allowbreak{}Weak\_\allowbreak{}Bianchi.md} proves

\[
 |\nabla v|^2\to2m_*,\quad |C|^2\to4m_*,\quad
 |c|^2\to1\quad\text{in measure},
 \qquad c=C/(2|v|).
\tag{15}
\]

For completeness, the new covariance test is \(T(q)=q/(1+q)\),
\(q=|\nabla v|^2\). Since \(\Delta|v|^2=4|v|^2-2q+o_{L^2}(1)\),
\(d|v|^2\to0\) in \(L^2\), and \(|dT(q)|\le C|\nabla^2v|\), testing
yields

\[
 \fint\frac{(q-2m_*)^2}{1+q}\longrightarrow0.
\]

Define the actual polynomial section

\[
 \sigma=\Phi\wedge t(v)\in\Gamma(L^5).
\]

The bundle power follows from \(t(v)\in S^+\otimes L^4\) and
\(\det S^+=L\). In the surviving first-jet configuration,

\[
 t(v)=(0,2|v|^3c),\qquad
 |\sigma|^2\to4\rho_*m_*^3>0.
\]

Moreover

\[
 |\sigma|\le C|\nabla v|,
\quad
 |\nabla\sigma|\le C\bigl(
 |\nabla\Phi||\nabla v|+|\nabla v|^2+|\nabla^2v|\bigr),
\]

so \(\|\sigma\|_4+\|\nabla\sigma\|_2\le C_B\). In particular
\(|\sigma|\) tends to its positive constant strongly in \(L^2\).

The three hypotheses of the line lemma now hold:

\begin{itemize}
\tightlist
\item
  \(\operatorname{Ric}=-g+o_{L^2}(1)\), so its normalized square is
  uniformly bounded.
\item
  The curvature of \(L^5\) is \(5F_A\). Its self-dual part is bounded
  pointwise by the actual spinor equation; its anti-self-dual part has
  the previously proved bounded same-field fit plus an \(L^2\)-small
  error.
\item
  The section \(\sigma\) has the amplitude and first-derivative bounds
  just proved.
\end{itemize}

The lemma forces

\[
 \frac{c_1(L^5)^2[M_j]}{V_j}\to0.
\]

But the actual endpoint characteristic budgets give

\[
 \frac{c_1(L^5)^2[M_j]}{V_j}
 =25\frac{a^2[M_j]}{V_j}
 \longrightarrow\frac{25}{6\pi^2}>0.
\tag{16}
\]

This is a contradiction. Every fixed finite-density zero-end sequence is
therefore excluded.

\hypertarget{returning-to-all-parameters-and-the-original-metric}{%
\subsection{4. Returning to all parameters and the original
metric}\label{returning-to-all-parameters-and-the-original-metric}}

The existing finite-parameter estimate excludes a sequence with
vanishing ratio in (1) and bounded \(\beta\); hence failure of (1) at a
fixed \(B\) would force \(\beta\to\infty\) and \(\kappa_\beta\to0\). The
exact Yamabe transfer of the manuscript gives

\[
 \widehat\kappa_\beta\le
 \frac{\kappa_\beta}{1-\sqrt{\kappa_\beta}-\kappa_\beta/16}
\]

on its small-defect branch. Thus such a failure on the original metrics
supplies precisely a scalar-minus-four sequence of the form (14), which
has just been excluded. This proves (1).

\hypertarget{exact-remaining-boundary-of-this-result}{%
\subsection{5. Exact remaining boundary of this
result}\label{exact-remaining-boundary-of-this-result}}

The theorem now applies to every prescribed finite \(B\), rather than
only \(B<39/4+\delta_*\). It completes the topology-density-controlled
version of O10 stated in the original research directive.

The proof uses \(B\) to bound \(\nabla^2v\), \(\nabla v\) in \(L^4\),
and ultimately \(\nabla\sigma\) in \(L^2\). It does not provide these
bounds uniformly if \(\chi/A\to\infty\). Consequently it does not prove
an absolute \(c>0\) for the full unrestricted O10 class. Such a
statement requires a further argument excluding simultaneous unbounded
opposite-Weyl energy and derivative energy, or a genuine counterexample
satisfying the original actual hypotheses.

\hypertarget{where-an-unbounded-opposite-weyl-budget-would-enter}{%
\subsubsection{5.1 Where an unbounded opposite-Weyl budget would
enter}\label{where-an-unbounded-opposite-weyl-budget-would-enter}}

Write \(b_j=\fint|W^-|^2\). The present proof has three specific uses of
a uniform bound on \(b_j\):

\begin{enumerate}
\def\labelenumi{\arabic{enumi}.}
\tightlist
\item
  The derivative-interchange estimate on \(E_-=S^-\otimes L\) gives
  \(\fint(|\nabla^2v|^2+|\nabla v|^4)\le C(1+b_j)\). The negative-spin
  connection contains \(W^-\); the estimate does not discard that block.
\item
  The scalar mass-gradient and covariance tests use
  \(\|\Delta|v|^2\|_2\) and \(\|\nabla^2v\|_2\). They therefore also
  need the same bound before the positive constant modulus of \(\sigma\)
  is obtained. For example, \(\|da\|_2^2\le\|a-m_*\|_2\|\Delta a\|_2\)
  pairs a small amplitude error with a quantity growing with \(b_j\).
\item
  The first-derivative estimate for \(\sigma\) uses \(\nabla^2v\),
  \(\nabla v\otimes\nabla v\), and \(\nabla\Phi\otimes\nabla v\). Their
  \(L^2\) bounds again depend on \(b_j\). The rank-one lemma requires a
  uniform bound along the sequence.
\end{enumerate}

The direct mass--phase--Dirac reduction avoids the earlier
Nijenhuis-bracket adjoint estimate entirely, but it still uses the
mass-gradient conclusion in item 2. A future estimate only bounding a
selected component of \(\nabla\sigma\) would not by itself justify the
unrestricted conclusion; it must also preserve the constant-modulus
conclusion, or replace that step. No rate comparison between the
original endpoint error and \(b_j\) is imposed in unrestricted O10, so
none may be inserted when passing to a subsequence.

\hypertarget{proof-provenance-within-this-continuation}{%
\subsection{6. Proof provenance within this
continuation}\label{proof-provenance-within-this-continuation}}

The rank-one logarithmic-resolvent lemma was developed in parallel in
\texttt{O10\_\allowbreak{}Topological\_\allowbreak{}Escape.md} and independently checked in
deriving Section 2 here. That note proves the more general pairing
conclusion \(V^{-1}\int F_j\wedge\omega_j\to0\) for a closed two-form
\(\omega_j\) with a bounded model and an \(L^2\)-small remainder. Only
the special case \(\omega_j=F_j\) is needed above.

The constant-length derivative and \(L^5\)-section calculations are
detailed in \texttt{O10\_\allowbreak{}Weak\_\allowbreak{}Bianchi.md}; their scalar test,
complex-frame norm factors, bundle exponent and derivative bound
received a separate algebraic check. The inherited strong-current result
remains a prerequisite, rather than a conclusion silently reproved by
the rank-one lemma.
\clearpage
\appendix
\hypertarget{strong-current-prerequisite-for-the-endpoint-proof}{%
\section{Strong-current prerequisite for the endpoint
proof}\label{strong-current-prerequisite-for-the-endpoint-proof}}

Date: 2026-09-29. Source: the supplied 20260929 integrated manuscript
(source hash in the accompanying status file), especially
\texttt{sec:single-parameter}, \texttt{sec:endpoint-covariance},
\texttt{sec:phase-curvature}, \texttt{sec:torsion-prolongation}, and
\texttt{sec:quartic-flux}.

\hypertarget{result-and-scope}{%
\subsection{Result and scope}\label{result-and-scope}}

There is a new global estimate available before attempting any further
manipulation of the final quartic flux identity. The actual
constant-weight positive spinor has uniformly integrable fourth-power
first derivative along every hypothetical O10 endpoint sequence. This
upgrades the cubic torsion Dirac equation, its length, the original
current derivative, and the full first torsion projection to strong
critical statements. It uses no diameter, Sobolev, noncollapse,
pointwise Ricci bound, or Euler-density upper bound.

This is an inherited prerequisite, reproduced to make the strengthened
endpoint proof independently readable. The present supplement
subsequently uses a bounded auxiliary section to obtain additional
global topological control; no fourth-power derivative estimate for the
original torsion is silently assumed.

All averages below use the actual total volume, even when the manifold
varies. Thus every vanishing squared norm is an \texttt{o(V)} statement,
not an unnormalized smallness assumption.

\hypertarget{a-slow-gradient-truncation-lemma}{%
\subsection{1. A slow gradient truncation
lemma}\label{a-slow-gradient-truncation-lemma}}

Let \((M_j,g_j)\) be closed Riemannian manifolds of fixed dimension and
\(E_j\) Hermitian bundles of fixed rank, with unitary connections. Let
\(u_j\) satisfy

\[
 |u_j|\le C_0,\qquad \nabla^*\nabla u_j=f_j(|u_j|^2)u_j,
 \qquad |f_j|+|f'_j|\le C_0.
\]

Suppose

\[
 \operatorname{Ric}_{g_j}=R_{0,j}+e_j,
 \quad F_j=F_{0,j}+E_j,
 \quad |R_{0,j}|+|F_{0,j}|\le C_0,
\]

where the symbol \(E_j\) on the right denotes a curvature error, and
assume

\[
 \mu_j:=\fint(|e_j|^2+|E_j|^2)\longrightarrow0,
 \qquad \sup_j\mu_j\le C_0.
\]

No independent Hessian or fourth-power gradient budget is needed: these
follow from the displayed hypotheses by the unweighted version of the
argument below.

There is \(\epsilon>0\), depending only on the displayed constants,
dimension, and rank, such that, with

\[
 q=|\nabla u|^2,\quad H=\nabla^2u,\quad
 w_L=\left(\frac q{L^2+q}\right)^\epsilon,
 \qquad L\ge1,
\]

one has

\[
 \boxed{\fint w_L(|H|^2+q^2)
              \le C L^{-2\epsilon}+C\mu_j.}\tag{A.1}
\]

Consequently

\[
 \boxed{\lim_{L\to\infty}\limsup_j
       \fint_{\{|\nabla u_j|\ge L\}}|\nabla u_j|^4=0.}\tag{A.2}
\]

This establishes uniform integrability of \(|\nabla u_j|^4\) relative to
normalized volume. Formula (1) also controls \(|H|^2\) on the
large-gradient set. It does \textbf{not}, by itself, prove uniform
integrability of \(|H|^2\) on arbitrary small-volume sets.

\hypertarget{proof}{%
\subsubsection{Proof}\label{proof}}

Use the smooth regularization

\[
 w_{L,\delta}=\left(\frac{q+\delta}{L^2+q+\delta}\right)^\epsilon,
 \quad \delta>0.
\]

It obeys

\[
 |dw_{L,\delta}|\,|\nabla u|\le2\epsilon w_{L,\delta}|H|.
\tag{A.3}
\]

Put \(P_i=\nabla_i u\), \(H_{ij}=\nabla_iP_j\), and \(I=\int w|H|^2\),
\(J=\int wq^2\), \(A_w=\int w(1+q)\). Antisymmetry gives the exact
algebraic identity

\[
 I=\int w\operatorname{Re}\langle H_{ij},H_{ji}\rangle
       +\frac12\int w\sum_{i,j}|F_{ij}u|^2.
\tag{A.4}
\]

In the first term integrate the derivative indexed by \(j\) in
\(H_{ji}\). Then commute the two derivatives in
\(\nabla_jH_{ij}=\nabla_j\nabla_iP_j\). The commutator on the
covector-valued section \(P\) contains only
\(\operatorname{Ric}*P+F*P\), and \(\sum_j\nabla_jP_j=-f(|u|^2)u\). In
particular \textbf{no derivative of either curvature occurs}. The
derivative of \(f(|u|^2)u\) is bounded by \(C|P|\). The derivative of
the scalar weight contributes at most \(2\epsilon I\), by (3). Thus,
after taking \(\epsilon\) sufficiently small,

\[
 I\le C A_w+C\int w(|e|+|E|)q+C\int (|e|^2+|E|^2).
\tag{A.5}
\]

The term involving \(F^2\) in (4) was bounded using \(F=F_0+E\) and
bounded \(F_0\). No derivative of the bounded model is required.

For completeness, a weighted bounded-section interpolation follows by
integrating the divergence of \(wq\operatorname{Re}\langle u,P\rangle\).
The derivative of \(q\) contributes at most \(C\int wq|H|\); the
derivative of \(w\) has the same bound by (3); the trace Hessian term is
bounded by \(C\int wq\). Therefore

\[
 J\le C\int wq+C\sqrt{IJ},\qquad
 J\le C_1(I+A_w).\tag{A.6}
\]

Since \(0\le w\le1\), for any fixed \(a>0\),

\[
 \int w(|e|+|E|)q
 \le aJ+C_a\int(|e|^2+|E|^2).
\tag{A.7}
\]

Choose \(a\) after \(C_1\) so that (5)--(7) absorb the resulting
multiple of \(I\). It follows that

\[
 I+J\le C A_w+C\mu_jV_j.
\tag{A.8}
\]

Before estimating the weighted right side, take \(w\equiv1\) in
(4)--(8). Testing the section equation by \(u\) gives \[
\int q=\int f(|u|^2)|u|^2\le CV.
\] The same absorption then proves \[
\fint(|H|^2+q^2)\le C(1+\mu_j)\le C.
\] Thus the energy budget used in the following step is a consequence,
not an extra hypothesis.

For \(0<\epsilon\le1\),

\[
 w_{L,\delta}\le L^{-2\epsilon}(q+\delta)^\epsilon.
\]

The just-proved normalized \(q^2\) bound controls the average of
\((q+\delta)^\epsilon(1+q)\), uniformly for \(0<\delta\le1\). Divide (8)
by \(V_j\), then let \(\delta\downarrow0\), obtaining (1). On
\(q\ge L^2\), \(w_L\ge2^{-\epsilon}\), so (2) follows. Every calculation
is on a closed manifold with a globally smooth regularized scalar
weight.

\hypertarget{apply-the-lemma-to-the-hypothetical-o10-endpoint}{%
\subsection{2. Apply the lemma to the hypothetical O10
endpoint}\label{apply-the-lemma-to-the-hypothetical-o10-endpoint}}

Use the manuscript's exact Yamabe reduction and a single ordinary
solution rescaled to

\[
 s=-4,\quad D_A\Phi=0,\quad F_A^+=2i\sigma(\Phi),
 \quad \rho=|\Phi|^2\le2.
\]

Set \(P=\nabla_A\Phi\), \(H=\nabla_AP\), \(\zeta=\langle\Phi,P\rangle\),
\(\tau=t_\Phi\), \(m=|\tau|^2\), \(\rho_*=4/3\), \(m_*=128/81\), and
\(\eta=r+D_\beta/V+\beta^{-2}\to0\).

The actual semilinear equation is

\[
 \nabla_A^*\nabla_A\Phi=(1-\rho/2)\Phi.\tag{A.9}
\]

The latest bounded ASD fit supplies the lemma's full
connection-curvature hypothesis. The old Hessian budget is consistent
with, and can also be recovered from, the unweighted part of the lemma.
With

\[
 f=-iF_A^-,\quad Q=-iq_-(\tau),\quad
 S=m_*+2\rho^2|P|^2,\quad
 K=-\frac98\frac{\rho^3Q}{S},
\]

the manuscript proves

\[
 |K|\le C,\qquad \fint|f-K|^2\le C\sqrt\eta.
\tag{A.10}
\]

The traceless curvature of \(S^+\) contains \(W^++s/12\) and the
off-diagonal Ricci block; it contains \textbf{no \(W^-\)}. Its central
part is \(F_A/2\). Retain the scalar block, the bounded actual self-dual
determinant curvature, and the bounded ASD model \(iK\) as \(F_0\). The
remaining curvature has

\[
 \fint|F^{S^+}-F_0|^2\le C\sqrt\eta,
 \quad |F_0|\le C,
 \quad \fint|\operatorname{Ric}+g|^2\le\eta.
\tag{A.11}
\]

Here \(W^+\to0\) follows from the same-solution parameter estimate,
including its bounded axial model of size \(O(\beta^{-1})\). Applying
(1) gives

\[
 \boxed{\fint w_L(|H|^2+|P|^4)
         \le C L^{-2\epsilon}+C\sqrt\eta.}\tag{A.12}
\]

There is no hypothesis on the opposite Weyl curvature and hence no use
of \(\chi/A\le B\) in this result.

\hypertarget{strong-critical-consequences}{%
\subsection{3. Strong critical
consequences}\label{strong-critical-consequences}}

\hypertarget{the-mixed-spinor-product-and-torsion-dirac-residual}{%
\subsubsection{3.1 The mixed spinor product and torsion Dirac
residual}\label{the-mixed-spinor-product-and-torsion-dirac-residual}}

The old identity gives \(\fint|\zeta|^2\le C\eta\), while
\(|\zeta|\le\sqrt2|P|\). Splitting into \(|P|\le L\) and its complement
in (12) yields

\[
 \fint|\zeta|^2|P|^2
 \le C\{L^2\eta+L^{-2\epsilon}+\sqrt\eta\}\to0.
\tag{A.13}
\]

For example choosing \(L=\eta^{-1/(2+2\epsilon)}\) gives a bound by
\(C\eta^{\epsilon/(1+\epsilon)}+C\sqrt\eta\). Also
\(\fint|\zeta|^4\to0\). The exact cubic equation, with axial
cancellation retained, now implies

\[
 \boxed{\|D^-_{3A}\tau\|_{L^2_{\rm avg}}\to0.}\tag{A.14}
\]

This is stronger than the previous \(L^q\), \(q<2\), conclusion.

\hypertarget{torsion-length-and-positive-spinor-energy-density}{%
\subsubsection{3.2 Torsion length and positive-spinor energy
density}\label{torsion-length-and-positive-spinor-energy-density}}

The existing torsion-length concentration gives \(m\to m_*\) in
normalized measure. Since \(m\le2\rho^2|P|^2\le8|P|^2\), (12) makes
\(m^2\) uniformly integrable. Vitali on the probability spaces
\((M_j,dV_j/V_j)\) gives

\[
 \boxed{\|m-m_*\|_{L^2_{\rm avg}}\to0.}\tag{A.15}
\]

The exact algebraic defect satisfies

\[
 \nu=2\rho^2|P|^2-m,\quad
 0\le\nu\le4\rho|\zeta|^2,
\]

so \(\nu\to0\) in normalized \(L^2\). On \(\rho\ge c>0\), use (15) and
\(\rho\to\rho_*\) in measure; on \(\rho<c<\rho_*\), its relative volume
tends to zero and (12) controls \(|P|^4\). Consequently

\[
 \boxed{\big\||P|^2-4/9\big\|_{L^2_{\rm avg}}\to0.}\tag{A.16}
\]

The scalar equation \(\Delta\rho=-2|P|^2+2\rho-\rho^2\) therefore gives

\[
 \|\Delta\rho\|_{L^2_{\rm avg}}\to0.\tag{A.17}
\]

\hypertarget{the-full-current-derivative}{%
\subsubsection{3.3 The full current
derivative}\label{the-full-current-derivative}}

The actual polynomial curvature equation is

\[
 \rho^3 f+2Q=2\rho^2(d\gamma)^-+\mathcal R,
 \quad \gamma=\operatorname{Im}\zeta,
 \quad |\mathcal R|\le
 C(\rho^3|\operatorname{Ric}^{\circ}|+
              \rho^{3/2}|\zeta||P|).
\tag{A.18}
\]

Real/imaginary conventions may reverse both signs of the derivative
term; only its norm is used here. By (13), \(\mathcal R\to0\) in
normalized \(L^2\). Furthermore

\[
 \rho^3 f+2Q
 =\rho^3(f-K)+\left(2-\frac98\frac{\rho^6}{S}\right)Q.
\tag{A.19}
\]

The first term is small by (10). The coefficient in the second is
uniformly bounded, converges to zero in measure, and \(|Q|^2=m^2/8\) is
uniformly integrable by (15). Thus the second term is small in
normalized \(L^2\). Equation (18) proves \(\rho^2(d\gamma)^-\to0\) in
that norm.

This weighted assertion becomes unweighted without assuming a pointwise
lower bound for \(\rho\). Indeed

\[
 |d\zeta|\le C(|P|^2+\rho|F^{S^+}|).
\tag{A.20}
\]

The squares on the right are uniformly integrable, respectively by (12)
and the bounded-model fit (11). The set \(\{\rho<c\}\), for any fixed
\(0<c<\rho_*\), has vanishing relative volume. On its complement divide
by \(\rho^2\); on that set use (20). Hence

\[
 \|(d\gamma)^-\|_{L^2_{\rm avg}}\to0.
\tag{A.21}
\]

The real part \(\operatorname{Re}\zeta=d\rho/2\) is exact, and
\(\gamma\) is a globally smooth one-form with \(\delta\gamma=0\). Stokes
gives exactly

\[
 \int|d\gamma^+|^2=\int|d\gamma^-|^2,
\]

so \(d\gamma\to0\) in normalized \(L^2\). The one-form and scalar
Bochner formulas yield

\[
 \int|\nabla\gamma|^2
 =\int|d\gamma|^2-\int\operatorname{Ric}(\gamma,\gamma),
\]

\[
 \int|\nabla^2\rho|^2
 =\int(\Delta\rho)^2-\int\operatorname{Ric}(d\rho,d\rho).
\]

The \(-g\) terms are small by the original current bound. The
Ricci-error pairings are bounded by its \(L^2_{\rm avg}\) norm times the
squared \(L^4_{\rm avg}\) norm of \(\gamma\) or \(d\rho\), which tend to
zero by (13). Therefore

\[
 \boxed{\|\nabla\zeta\|_{L^2_{\rm avg}}\to0.}\tag{A.22}
\]

This proves the conclusion sought from the manuscript's current--Hessian
correlation \(\mathcal B\) \textbf{without adding \(\mathcal B\to0\) as
a hypothesis} and without claiming that \(\mathcal B\) itself has been
estimated.

\hypertarget{the-full-first-torsion-projection}{%
\subsubsection{3.4 The full first torsion
projection}\label{the-full-first-torsion-projection}}

Use the already proved actual identity

\[
 N=\sqrt{\rho/2}\,\mathcal L(\nabla\zeta)+\mathcal R_N,
\quad
 |\mathcal R_N|\le C\{\rho^{3/2}(
 |\operatorname{Ric}^{\circ}|+|W^+-W_0|)+|\zeta||P|\}.
\]

Equations (13), (22), and the curvature errors give

\[
 \boxed{\|N\|_{L^2_{\rm avg}}\to0.}\tag{A.23}
\]

This is the complete displayed first torsion projection, not merely its
Dirac trace.

\hypertarget{the-remaining-torsion-gradient-has-a-fixed-nonzero-mean}{%
\subsubsection{3.5 The remaining torsion gradient has a fixed nonzero
mean}\label{the-remaining-torsion-gradient-has-a-fixed-nonzero-mean}}

The unweighted twisted Lichnerowicz identity is now legitimate at the
limiting strength:

\[
 \fint|\nabla_{3A}\tau|^2
 =\fint|D^-_{3A}\tau|^2+\fint m
       -\frac32\fint\langle c(F_A^-)\tau,\tau\rangle.
\]

Replace \(f=-iF_A^-\) by its bounded fit \(K\) using Cauchy--Schwarz
with the uniformly bounded \(L^2\) norm of \(m\). The Clifford identity
and (15) then give

\[
 \frac32\fint\langle c(F_A^-)\tau,\tau\rangle
 =-\frac{27}{32}\fint\frac{\rho^3m^2}{S}
 \longrightarrow-m_*.
\]

The bounded coefficient \(\rho^3/S\), convergence in measure, and
uniform integrability of \(m^2\) justify this replacement. Consequently

\[
 \boxed{\fint|\nabla_{3A}\tau|^2\longrightarrow
              2m_* =256/81.}\tag{A.25}
\]

Thus the strengthened first projection does not imply that the torsion
is parallel: its surviving derivative energy is nonzero and sits in the
complementary type component.
\clearpage
\hypertarget{direct-bounded-resolvent-and-derivative-estimates}{%
\section{Direct bounded resolvent and derivative
estimates}\label{direct-bounded-resolvent-and-derivative-estimates}}

This appendix adopts and verifies the direct construction in the newly
supplied O10 resolvent attachment. It is used as a proof simplification,
not counted as a second new endpoint theorem.

\hypertarget{direct-resolvent-complete-energy-estimate}{%
\subsection{2. Direct resolvent: complete energy
estimate}\label{direct-resolvent-complete-energy-estimate}}

Use normalized volume throughout. Let \texttt{t} be the actual cubic
negative spinor, \texttt{L\ =\ nabla*\ nabla}, and \texttt{D} its
twisted Dirac operator. The inherited endpoint conclusions are

\[
 \|Dt\|_2\to0,\qquad |t|^2\to m_* =128/81\text{ in }L^2,
 \qquad (L-2)t=D^*Dt+(\gamma-1)t-et,
\]

where \texttt{gamma} is uniformly bounded and converges to 1 in measure,
while \texttt{e\ -\textgreater{}\ 0} in normalized \texttt{L2}. Thus
\texttt{gamma-1\ -\textgreater{}\ 0} in every finite normalized
\texttt{Lp}, and \texttt{t} is bounded in normalized \texttt{L4}.

Choose a smooth radial bundle map \texttt{T} equal to the identity for
\texttt{\textbar{}t\textbar{}\ \textless{}=\ 2} and bounded in norm by
4. Set \texttt{z=T(t)}. Since \texttt{m\_*\textless{}4} and
\texttt{\textbar{}t\textbar{}²\ -\textgreater{}\ m\_*} strongly in
\texttt{L2},

\[
 \|z-t\|_2\to0.
\]

No derivative of this cutoff is required. Define

\[
 (4+L)u=6z.
\tag{B.2.1}
\]

The connection-Laplacian maximum principle gives

\[
 |u|\le6,\qquad |Lu|\le48.
\tag{B.2.2}
\]

Put \texttt{w=u-t} and \texttt{d=Dt}. Then

\[
 (4+L)w=6(z-t)-D^*d-(\gamma-1)t+et.
\]

Pairing with \texttt{w} gives

\[
 4\|w\|_2^2+\|\nabla w\|_2^2
 =6\langle z-t,w\rangle-\langle d,Dw\rangle
 -\langle(\gamma-1)t,w\rangle+\langle et,w\rangle.
\]

The first two terms are absorbed by Young's inequality, using
\texttt{\textbar{}Dw\textbar{}\ \textless{}=\ 2\textbar{}nabla\ w\textbar{}}.
The remaining two terms satisfy

\[
 |\langle(\gamma-1)t,w\rangle|
 \le\|\gamma-1\|_2\|t\|_4\|w\|_4=o(1),
\]

\[
 |\langle et,w\rangle|
 \le\|e\|_2\|t\|_4\|w\|_4=o(1).
\]

Here
\texttt{\textbar{}\textbar{}w\textbar{}\textbar{}4\ \textless{}=\ \textbar{}\textbar{}u\textbar{}\textbar{}infinity\ +\ \textbar{}\textbar{}t\textbar{}\textbar{}4}
follows directly on the normalized probability space. It is not inferred
from a Sobolev estimate. Therefore

\[
 \boxed{\|u-t\|_{H^1_{\rm av}}\to0.}
\tag{B.2.3}
\]

The same direct resolvent also supplies the residual needed in the
stronger \texttt{39/4} argument:

\[
 \boxed{(L-2)u=6(z-u)\to0\text{ in }L^2_{\rm av},\qquad Du\to0\text{ in }L^2_{\rm av}.}
\tag{B.2.4}
\]

Thus all uses of the old bounded resolvent section remain valid after
this substitution. In particular
\texttt{\textbar{}u\textbar{}²\ -\textgreater{}\ m\_*} in normalized
\texttt{L2}, because \texttt{u} is bounded and converges in measure to
the same amplitude as \texttt{t}.

\hypertarget{why-the-weighted-hessian-estimate-does-not-hide-extra-regularity}{%
\subsection{3. Why the weighted Hessian estimate does not hide extra
regularity}\label{why-the-weighted-hessian-estimate-does-not-hide-extra-regularity}}

Write \texttt{P\_\allowbreak{}i=nabla\_\allowbreak{}i\ u}, \texttt{H\_\allowbreak{}ij=nabla\_\allowbreak{}i\ P\_\allowbreak{}j},
\texttt{q=\textbar{}P\textbar{}²}, and \texttt{h=Lu}. The exact
antisymmetry identity is

\[
 \int w|H|^2
 =\int w\operatorname{Re}\langle H_{ij},H_{ji}\rangle
 +\frac12\int w\sum_{i,j}|F_{ij}u|^2.
\]

In the first term, integrate the derivative with index \texttt{j}, then
commute \texttt{nabla\_\allowbreak{}j\ nabla\_\allowbreak{}i\ P\_\allowbreak{}j}. The contracted covector
curvature is Ricci; the other curvature acts on the bundle slot. The
only derivative of a traced quantity is \texttt{nabla\ h}. Integrating
that term back once gives

\[
 \int w\operatorname{Re}\langle P,\nabla h\rangle
 =\int w|h|^2-\int\operatorname{Re}\langle dw\otimes P,h\rangle.
\]

No \texttt{nabla\ F}, \texttt{nabla\ Ric}, or estimate on
\texttt{nabla\ h} is introduced. For

\[
 w=\left(\frac{q+\delta}{L_0^2+q+\delta}\right)^\epsilon,
 \qquad |dw|\sqrt q\le2\epsilon w|H|,
\]

the weight terms are absorbed for a fixed sufficiently small epsilon.
The bounded-section interpolation from
\texttt{div(wq\ Re\textless{}u,nabla\ u\textgreater{})} gives

\[
 J=\int wq^2\le C_1\left(I+\int wq\right),
 \qquad I=\int w|H|^2.
\]

Split the bundle curvature at an arbitrary amplitude \texttt{T}. Its low
part contributes \texttt{C\_\allowbreak{}T\ int\ w(1+q)}. Young's inequality treats
the tail-curvature times \texttt{q} term with a small multiple of
\texttt{J}; choose that small coefficient after fixing \texttt{C1}. The
coefficient multiplying the tail's squared norm is then independent of
\texttt{T} and of the total curvature-energy bound. Consequently

\[
 \fint w(|H|^2+q^2)
 \le C_TL_0^{-2\epsilon}
 +C\fint_{|F|>T}|F|^2+C\fint|\operatorname{Ric}+g|^2.
\tag{B.3.1}
\]

The unweighted calculation first supplies the finite \texttt{L2} Hessian
and \texttt{L4} gradient budgets. Their constants can depend on the
total curvature-energy bound; this dependence is permitted in
\texttt{C\_\allowbreak{}T} and disappears when \texttt{L0\ -\textgreater{}\ infinity}
with \texttt{T} fixed.

For any sets \texttt{Z\_\allowbreak{}j} of vanishing relative volume, let

\[
 \Gamma_F=\lim_{T\to\infty}\limsup_j\fint_{|F_j|>T}|F_j|^2.
\]

First take \texttt{j\ -\textgreater{}\ infinity}, then
\texttt{L0\ -\textgreater{}\ infinity}, then
\texttt{T\ -\textgreater{}\ infinity} in (3.1). This proves

\[
 \limsup_j\fint_{Z_j}|\nabla u_j|^4\le C\Gamma_F.
\tag{B.3.2}
\]

This is a tail estimate, not an assertion that the Hessian squares are
uniformly integrable on arbitrary small sets.
\clearpage
\hypertarget{o10-the-regularized-derivative-has-asymptotically-constant-length}{%
\section{O10: the regularized derivative has asymptotically constant
length}\label{o10-the-regularized-derivative-has-asymptotically-constant-length}}

Date: 2026-09-29. This note continues the existing actual O10 endpoint
argument for a sequence with a fixed finite bound on the normalized
opposite Weyl energy. It proves a further consequence of the bounded
resolvent construction and identifies a genuine next coefficient bundle.
It does \textbf{not} prove unrestricted O10.

\hypertarget{inputs-already-proved-in-continuation-b}{%
\subsection{1. Inputs already proved in continuation
B}\label{inputs-already-proved-in-continuation-b}}

All norms and integrals below are normalized by volume. Let

\[
 v\in\Gamma(E_-),\qquad E_-=S^-\otimes L,\qquad
 a=|v|^2,\quad q=|\nabla v|^2,\quad m_*=128/81.
\]

For each fixed bound on \(\fint |W^-|^2\), the preceding resolvent
argument gives

\[
 \|v\|_\infty\le C,\qquad
 \|\nabla v\|_4+\|\nabla^2v\|_2\le C_B,
 \qquad a\to m_*\text{ in }L^p\ (p<\infty),
\]

\[
 e=\nabla^*\nabla v-2v\to0\text{ in }L^2,
 \qquad da\to0\text{ in }L^2.
\tag{C.1}
\]

The last assertion was obtained from the scalar identity for \(a\), its
strong amplitude convergence, and the uniform \(L^2\) bound on \(q\).
The strong-current prerequisite gives \(N_v=(\nabla v)^{1,0}_{J_+}\to0\)
in normalized \(L^2\). Together with (1) and \(D^-_{3A}v\to0\), this
gives, in compatible frames adapted to the actual positive spinor and to
\(v\),

\[
 \nabla_{Z_1}v\to0,\quad \nabla_{Z_2}v\to0,
 \quad \nabla_{\bar Z_1}v\to0,
 \quad (\nabla_{\bar Z_2}v)_1\to0
 \qquad\text{in }L^2.
\tag{C.2}
\]

Statements in adapted frames are understood invariantly on fixed
positive-amplitude sets, with the exceptional sets removed using the
established \(L^4\) derivative bound. No smooth limiting manifold or
noncollapse is assumed.

Here is the direct proof of (2), avoiding a Nijenhuis bracket or a
second-derivative adjoint test. Write \[
 A_i=\nabla_{Z_i}v,\quad B=\nabla_{\bar Z_1}v,
 \quad C=\nabla_{\bar Z_2}v,\quad
 j_v=\operatorname{Im}\langle v,\nabla v\rangle.
\] The exact elementary phase identity \[
 |j_v+\tfrac12J_+da|^2\le2a|N_v|^2
\] shows that \(j_v\to0\) in \(L^2\), since \(v\) is bounded and
\(da\to0\). In a compatible frame with \(v=\sqrt a\,e_1\), every first
spinor component obeys \[
 (\nabla_Xv)_1=\frac{da(X)/2+i j_v(X)}{\sqrt a}
\] for real \(X\), extended complex linearly. On a fixed set
\(a\ge c>0\), all these components therefore tend to zero in \(L^2\).
The sets where \(a<c<m_*\), or where the actual positive amplitude
\(\rho<c'<\rho_*\), have relative volume tending to zero; there \[
 \fint_E|\nabla v|^2
 \le \mu(E)^{1/2}\bigl(\fint|\nabla v|^4\bigr)^{1/2}\to0.
\] Hence \(B_1,C_1\to0\) globally in the invariant cutoff sense. The
actual Clifford formula is \[
 D^-v=((A_1)_1-(A_2)_2,\ B_2+C_1),
\] up to an overall sign. It follows that \(B_2\to0\). Finally
\(A_1,A_2\to0\) are exactly the small \((1,0)\) projection supplied by
\(N_v\). This proves (2) using only mass, phase, and Dirac information.

\hypertarget{a-scalar-covariance-argument-gives-more-than-the-averaged-energy-limit}{%
\subsection{2. A scalar covariance argument gives more than the averaged
energy
limit}\label{a-scalar-covariance-argument-gives-more-than-the-averaged-energy-limit}}

Use the nonnegative scalar Laplacian. Equation (1) gives

\[
 \Delta a=4a-2q+2\operatorname{Re}\langle e,v\rangle.
\tag{C.3}
\]

Set \(q_*=2m_*\) and

\[
 T(x)=\frac{x}{1+x},\qquad x\ge0.
\]

The chain rule and the covariant Kato estimate imply

\[
 |dT(q)|\le \frac{2\sqrt q}{(1+q)^2}|\nabla^2v|
 \le C|\nabla^2v|.
\tag{C.4}
\]

Thus \(T(q)-T(q_*)\) is a bounded test function with a uniformly bounded
normalized \(H^1\) seminorm, even though \(q\) itself need not have an
\(L^2\) gradient. Testing (3) with this function gives

\[
\begin{aligned}
 2\fint(q-q_*)[T(q)-T(q_*)]
 ={}&4\fint(a-m_*)[T(q)-T(q_*)]\\
 &+2\fint\operatorname{Re}\langle e,v\rangle[T(q)-T(q_*)]\\
 &-\fint\langle da,dT(q)\rangle\longrightarrow0.
\end{aligned}
\tag{C.5}
\]

Every term on the right tends to zero by (1) and (4). The integrand on
the left is nonnegative and is bounded below by a positive number on
every set \(|q-q_*|\ge\epsilon\). Consequently

\[
 \boxed{|\nabla v|^2\longrightarrow2m_*\quad\text{in normalized measure}.}
\tag{C.6}
\]

Since \(q\) is uniformly bounded in normalized \(L^2\), it follows that

\[
 \boxed{|\nabla v|^2\longrightarrow2m_*
       \quad\text{strongly in normalized }L^p\text{ for every }1\le p<2.}
\tag{C.7}
\]

This strengthens the preceding result \(\fint|\nabla v|^2\to2m_*\). It
does \textbf{not} establish strong \(L^2\) convergence of \(q\),
equivalently uniform integrability of \(|\nabla v|^4\). Critical
concentration remains possible under these estimates.

\hypertarget{the-remaining-first-derivative-coefficient-has-nonzero-constant-modulus}{%
\subsection{3. The remaining first-derivative coefficient has nonzero
constant
modulus}\label{the-remaining-first-derivative-coefficient-has-nonzero-constant-modulus}}

Let \(C=(\nabla_{\bar Z_2}v)_2\). Because \(Z_i=e_{2i-1}-ie_{2i}\), the
full derivative norm is

\[
 |\nabla v|^2=\tfrac12\bigl(
 |\nabla_{Z_1}v|^2+|\nabla_{\bar Z_1}v|^2+
 |\nabla_{Z_2}v|^2+|\nabla_{\bar Z_2}v|^2\bigr).
\]

Equations (2) and (7) therefore imply

\[
 \boxed{|C|^2\longrightarrow4m_*\quad\text{in measure and in }L^1.}
\tag{C.8}
\]

Write the surviving negative-spinor off-diagonal connection as

\[
 b=(0,0,c,-ic),\qquad c=C/(2\sqrt a),
\]

on fixed positive-amplitude sets. Then

\[
 \boxed{|c|^2\longrightarrow1\quad\text{in measure}.}
\tag{C.9}
\]

The conclusion concerns the invariant modulus. It does not provide a
global choice of the phase of \(c\).

\hypertarget{a-genuine-line-bundle-coefficient-for-the-next-prolongation}{%
\subsection{4. A genuine line-bundle coefficient for the next
prolongation}\label{a-genuine-line-bundle-coefficient-for-the-next-prolongation}}

The cubic construction is an algebraic first-order operation on any
spinor; it does not require that the spinor solve a monopole equation.
Since \(v\) is a negative spinor with determinant line \(L^3\), its
cubic torsion is a smooth positive spinor

\[
 t(v)\in\Gamma(S^+\otimes L^4).
\tag{C.10}
\]

Use the actual positive spinor \(\Phi\in\Gamma(S^+)\) to form the
determinant pairing

\[
 \boxed{\sigma=\Phi\wedge t(v)\in\Gamma(L^5).}
\tag{C.11}
\]

This is a globally defined section. The exponent is checked directly:
\(\det S^+=L\), and the second factor is twisted by \(L^4\).

In compatible frames \(\Phi=\sqrt\rho\,e_1\), \(v=\sqrt a\,e_1\), the
ideal remaining first jet gives

\[
 t(v)=(0,2a^{3/2}c),\qquad
 |\sigma|^2=4\rho a^3|c|^2.
\tag{C.12}
\]

The omitted components are bounded linear multiples of the projections
in (2); hence they tend to zero in \(L^2\). With \(\rho\to\rho_*=4/3\),
(9) gives

\[
 \boxed{|\sigma|^2\longrightarrow
 a_{\sigma,*}=4\rho_*m_*^3>0\quad\text{in measure}.}
\tag{C.13}
\]

The exact polynomial definition and boundedness of \(\Phi,v\) also give

\[
 |\sigma|\le C|\nabla v|,
\]

\[
 |\nabla\sigma|\le C\bigl(
 |\nabla\Phi|\,|\nabla v|+|\nabla v|^2+|\nabla^2v|\bigr).
\]

Therefore

\[
 \boxed{\|\sigma\|_4+\|\nabla\sigma\|_2\le C_B,
 \qquad |\sigma|^2\to a_{\sigma,*}\text{ strongly in }L^p\ (1\le p<2).}
\tag{C.14}
\]

Unlike the connection on \(E_-\), this line has curvature \(5F_A\),
already a bounded actual model plus an error tending to zero in
normalized \(L^2\). Its characteristic square has positive endpoint
density:

\[
 \frac{c_1(L^5)^2[M]}V=25\frac{a^2}{V}
 \longrightarrow\frac{25}{6\pi^2}.
\tag{C.15}
\]

These facts identify a concrete next coefficient, rather than an
abstract proposal to add a new field.

\hypertarget{what-has-not-yet-been-inferred}{%
\subsection{5. What has not yet been
inferred}\label{what-has-not-yet-been-inferred}}

A nearly constant nonzero modulus in measure and a bounded \(H^1\)
energy do not supply a nowhere-zero smooth section. In particular,
(13)--(15) alone are not a contradiction on varying geometries. A
separate geometry-free line-bundle compactness or actual differential
constraint would be required.

Nor has this note proved

\[
 D^+_{9A}t(v)\to0\quad\text{in }L^2.
\]

The Lichnerowicz potential along the surviving positive-spinor line is
negative, tending to \(-7\). That sign prevents a positive-potential
vanishing argument, but it does not itself prevent a
resolvent-and-characteristic-class argument if a sufficiently small
Dirac residual can be established. The latter remains unproved.

The exact anti-self-dual Einstein almost-Kähler calculation in
Armstrong's thesis, Section 3.3.3, constrains derivatives of the
analogous complex coefficient and then uses differential Bianchi
identities. Its smooth exact argument does not justify discarding
products of an \(L^2\)-small first projection, an \(L^4\)-bounded
derivative, and an \(L^2\)-bounded curvature in the present weak
setting. The covariance argument above closes a genuine amplitude step.
The rank-one argument in Section 6 bypasses this further weak
curvature/prolongation step for fixed finite density; it does not
establish the latter separately.

Primary reference checked: John Armstrong, thesis, Section 3.3.3, PDF
pp.~82--84,
https://johnarmstrongmaths.com/publications/resources/thesis.pdf . The
screenshot of PDF p.~83 is needed to distinguish the two different
spinor directions in the derivative constraints; plain-text extraction
loses one conjugation.

\hypertarget{integration-with-the-new-rank-one-lemma}{%
\subsection{6. Integration with the new rank-one
lemma}\label{integration-with-the-new-rank-one-lemma}}

The covariance identity used above has the explicit form

\[
 (q-q_*)[T(q)-T(q_*)]
 =\frac{(q-q_*)^2}{(1+q)(1+q_*)},
\]

so the proof also gives \(\fint (q-q_*)^2/(1+q)\to0\).

The additional line-bundle argument referred to in Section 5 has now
been supplied in \texttt{O10\_\allowbreak{}Topological\_\allowbreak{}Escape.md}. It uses the
actual normalized Ricci-energy bound, logarithmic selection of a
rough-Laplacian resolvent parameter, and rank-one Chern transgression.
Combining it with (13)--(15) produces the contradiction. The complete
combined proof is \texttt{O10\_\allowbreak{}All\_\allowbreak{}Finite\_\allowbreak{}Density.md}: O10 holds for
every fixed finite bound \(\chi/A\le B\), with a constant \(c(B)>0\).
The unrestricted uniform-density statement remains outside that proof.
\clearpage
\hypertarget{o10-a-rank-one-chern-obstruction-by-logarithmic-resolvent-selection}{%
\section{O10: a rank-one Chern obstruction by logarithmic resolvent
selection}\label{o10-a-rank-one-chern-obstruction-by-logarithmic-resolvent-selection}}

2026-09-29. This note proves an analytic lemma, then gives its exact
conditional application to the proposed third-torsion line section. The
lemma does not assume noncollapse, a Sobolev inequality, an
injectivity-radius bound, a Riemann-curvature bound, or convergence of
the manifolds. The application requires the separately derived
third-section amplitude and derivative estimates; those estimates are
not silently assumed to have been proved here.

\hypertarget{the-lemma}{%
\subsection{1. The lemma}\label{the-lemma}}

Let \((M_j,g_j)\) be closed oriented Riemannian four-manifolds, with
\(V_j=\operatorname{Vol}(M_j,g_j)\). All norms and integrals in the
proof below are normalized by \(V_j\). Let \(L_j\) be Hermitian line
bundles with unitary connections, and write their real curvature as
\(F_j=-iF_{\nabla_j}\). Assume

\[
\sup_j\left(\fint|\operatorname{Ric}_{g_j}|^2+\fint|F_j|^2\right)<\infty.
\tag{D.1.1}
\]

Suppose smooth sections \(\sigma_j\) satisfy

\[
|\sigma_j|\longrightarrow c>0\quad\text{in normalized measure},
\qquad \sup_j\fint|\nabla_j\sigma_j|^2<\infty.
\tag{D.1.2}
\]

Let \(\omega_j\) be closed real two-forms with a decomposition

\[
\omega_j=\omega_{0,j}+e_j,
\qquad \|\omega_{0,j}\|_\infty\le C,
\qquad \fint|e_j|^2\longrightarrow0.
\tag{D.1.3}
\]

Then

\[
\boxed{\fint F_j\wedge\omega_j\longrightarrow0.}
\tag{D.1.4}
\]

The notation in (1.4) means \(V_j^{-1}\int F_j\wedge\omega_j\); it is
not a pointwise inner-product convention. In particular, if \(F_j\)
itself admits (1.3), then

\[
\boxed{\frac{c_1(L_j)^2[M_j]}{V_j}\longrightarrow0.}
\tag{D.1.5}
\]

The uniform Ricci-energy assumption is essential to the argument. There
is no assertion here that a nearly unit section with bounded Dirichlet
energy alone rules out line-bundle topology on arbitrarily degenerating
metrics.

\hypertarget{uniformly-bounded-truncation}{%
\subsection{2. Uniformly bounded
truncation}\label{uniformly-bounded-truncation}}

Choose a smooth fiberwise radial map \(T\), equal to the identity on a
fixed neighborhood of the circle \(|z|=c\), uniformly bounded, and
globally Lipschitz. Put

\[
s_j=T(\sigma_j).
\]

This definition is invariant under unitary changes of frame. There are
constants \(K,E\), independent of \(j\), such that

\[
\|s_j\|_\infty\le K,\qquad
\fint|\nabla_j s_j|^2\le E,
\qquad d_j:=\fint (|s_j|-c)^2\longrightarrow0.
\tag{D.2.1}
\]

The last convergence follows from boundedness and convergence in
measure. No bound on the untruncated section's fourth power is needed.

\hypertarget{a-resolvent-parameter-selected-over-logarithmic-scales}{%
\subsection{3. A resolvent parameter selected over logarithmic
scales}\label{a-resolvent-parameter-selected-over-logarithmic-scales}}

Write \(\mathcal L_j=\nabla_j^*\nabla_j\). For \(\varepsilon>0\), solve
exactly

\[
u_{j,\varepsilon}=(I+\varepsilon\mathcal L_j)^{-1}s_j.
\tag{D.3.1}
\]

Each inverse exists on the individual compact manifold. The unitary
maximum principle gives \(\|u_{j,\varepsilon}\|_\infty\le K\). Spectral
calculus for the nonnegative self-adjoint operator \(\mathcal L_j\)
gives

\[
\|u_{j,\varepsilon}-s_j\|_2^2
\le\varepsilon\|\nabla_j s_j\|_2^2\le E\varepsilon,
\qquad
\|\nabla_j u_{j,\varepsilon}\|_2^2\le E,
\tag{D.3.2}
\]

and the exact integrated estimate

\[
\int_0^\infty\|\mathcal L_j u_{j,\varepsilon}\|_2^2\,d\varepsilon
=\|\nabla_j s_j\|_2^2\le E.
\tag{D.3.3}
\]

For an eigenvalue \(\lambda\ge0\), the identity behind (3.3) is
\(\int_0^\infty \lambda^2(1+\varepsilon\lambda)^{-2}\,d\varepsilon=\lambda\).

Set

\[
\delta_j=d_j+\|e_j\|_2^2+j^{-1}\longrightarrow0.
\]

For sufficiently large \(j\), choose \(\varepsilon_j\) in the interval

\[
\delta_j^{1/2}\le\varepsilon_j\le\delta_j^{1/4}
\tag{D.3.4}
\]

so that

\[
\varepsilon_j\|\mathcal L_j u_{j,\varepsilon_j}\|_2^2
\le\frac{4E}{|\log\delta_j|}\longrightarrow0.
\tag{D.3.5}
\]

Such a parameter exists by averaging (3.3) with respect to
\(d\varepsilon/\varepsilon\) on (3.4). Abbreviate
\(u_j=u_{j,\varepsilon_j}\) and \(K_j=\|\mathcal L_j u_j\|_2^2\). Thus

\[
\varepsilon_j\to0,\quad
\varepsilon_jK_j\to0,\quad
\frac{\|e_j\|_2^2}{\varepsilon_j}\le\sqrt{\delta_j}\to0.
\tag{D.3.6}
\]

This logarithmic selection is the decisive step. Choosing an arbitrary
smoothing parameter only gives \(K_j=O(\varepsilon_j^{-1})\), which
leaves a nonvanishing critical error in the argument below.

\hypertarget{derivative-free-hessian-estimate}{%
\subsection{4. Derivative-free Hessian
estimate}\label{derivative-free-hessian-estimate}}

For any bounded smooth section \(u\) of a unitary bundle on a closed
four-manifold, put

\[
I=\fint|\nabla^2u|^2,\qquad J=\fint|\nabla u|^4.
\]

Integrating the divergence of
\(|\nabla u|^2\operatorname{Re}\langle u,\nabla u\rangle\) gives

\[
J\le C\|u\|_\infty I^{1/2}J^{1/2},
\qquad J\le C\|u\|_\infty^2 I.
\tag{D.4.1}
\]

This is an integrated interpolation identity, not a Sobolev inequality.
Interchanging the two covariant derivative slots and integrating by
parts gives

\[
I\le \|\nabla^*\nabla u\|_2^2
+C\fint (|\operatorname{Ric}|+|F_{\nabla}|)|\nabla u|^2
+C\|u\|_\infty^2\|F_{\nabla}\|_2^2.
\tag{D.4.2}
\]

There is no derivative of \(F_{\nabla}\) in (4.2). One direct
verification is to apply the integrated covariant Hodge--Weitzenböck
identity to the bundle-valued one-form \(\nabla u\): its exterior
covariant derivative is \(F_{\nabla}u\), and its covariant
codifferential is \(\nabla^*\nabla u\). The curvature term on its
one-form index contains only Ricci, and the bundle term is algebraic in
\(F_{\nabla}\). Equivalently, use
\(\nabla_i\nabla_j u-\nabla_j\nabla_i u=F_{ij}u\), symmetrize the
Hessian square, and commute the remaining derivatives on the one-form
index. Cauchy--Schwarz, (4.1), and Young's inequality give, under (1.1),

\[
\boxed{\fint|\nabla u_j|^4+\fint|\nabla^2u_j|^2
\le C(1+K_j).}
\tag{D.4.3}
\]

Here \(C\) depends on the uniform bounds in (1.1), (2.1), and dimension,
and not on \(\varepsilon_j\), charts, or a Sobolev constant.

\hypertarget{small-quadratic-energy-on-the-amplitude-transition-region}{%
\subsection{5. Small quadratic energy on the amplitude-transition
region}\label{small-quadratic-energy-on-the-amplitude-transition-region}}

Fix \(0<a<b<c\), and let

\[
B_j=\{|u_j|\le b\}.
\]

By (2.1), (3.2), and (3.4),

\[
\mu_j(B_j)\le C(d_j+E\varepsilon_j)\le C\varepsilon_j,
\qquad \mu_j=dV_j/V_j.
\tag{D.5.1}
\]

Consequently (4.3) and (3.6) imply

\[
\fint_{B_j}|\nabla u_j|^2
\le \mu_j(B_j)^{1/2}
 \left(\fint|\nabla u_j|^4\right)^{1/2}
\le C\sqrt{\varepsilon_j(1+K_j)}\longrightarrow0.
\tag{D.5.2}
\]

At the same time,

\[
\|e_j\|_2
\left(\fint|\nabla u_j|^4\right)^{1/2}
\le C\sqrt{\|e_j\|_2^2(1+K_j)}\longrightarrow0,
\tag{D.5.3}
\]

because
\(\|e_j\|_2^2 K_j=(\|e_j\|_2^2/\varepsilon_j)(\varepsilon_jK_j)\to0\).

\hypertarget{rank-one-transgression}{%
\subsection{6. Rank-one transgression}\label{rank-one-transgression}}

On \(\{u_j\ne0\}\), the unit section \(u_j/|u_j|\) trivializes the line
bundle. Its real connection form \(\alpha_j\) satisfies

\[
d\alpha_j=F_j,\qquad
|\alpha_j|\le C|\nabla u_j|/|u_j|.
\]

Choose a fixed smooth cutoff \(\chi\), zero on \([0,a]\), one on
\([b,\infty)\), and put \(\chi_j=\chi(|u_j|)\). Then \(\chi_j\alpha_j\)
is a globally smooth real one-form. Since \(d\omega_j=0\),

\[
\int F_j\wedge\omega_j
=\int(1-\chi_j)F_j\wedge\omega_j
-\int d\chi_j\wedge\alpha_j\wedge\omega_j.
\tag{D.6.1}
\]

The first term tends to zero after division by \(V_j\): by (1.1),

\[
\left|\fint(1-\chi_j)F_j\wedge\omega_j\right|
\le C\|F_j\|_2\mu_j(B_j)^{1/2}
+\|F_j\|_2\|e_j\|_2\to0.
\tag{D.6.2}
\]

The second term is supported where \(a\le|u_j|\le b\), and there
\(|d\chi_j\wedge\alpha_j|\le C|\nabla u_j|^2\). Decomposing \(\omega_j\)
as in (1.3) gives

\[
\left|\fint d\chi_j\wedge\alpha_j\wedge\omega_j\right|
\le C\fint_{B_j}|\nabla u_j|^2
+C\|e_j\|_2\left(\fint|\nabla u_j|^4\right)^{1/2}\to0
\tag{D.6.3}
\]

by (5.2)--(5.3). Equations (6.1)--(6.3) prove the lemma. Taking
\(\omega_j=F_j\) and using \(c_1(L_j)=[F_j/(2\pi)]\) proves (1.5).

\hypertarget{conditional-application-to-the-third-torsion-line}{%
\subsection{7. Conditional application to the third torsion
line}\label{conditional-application-to-the-third-torsion-line}}

Suppose the weak third-torsion analysis produces, from a hypothetical
O10 zero sequence, a section

\[
\sigma_j\in\Gamma(L_j^{\otimes5})
\]

with

\[
|\sigma_j|\to c_\sigma>0\text{ in normalized measure},
\qquad \sup_j\fint|\nabla_{5A_j}\sigma_j|^2<\infty.
\tag{D.7.1}
\]

The original zero-endpoint assumptions give bounded normalized Ricci
energy, and the actual determinant-curvature fit gives

\[
-iF_{5A_j}=5(-iF_{A_j})=F_{0,j}+e_j,
\qquad |F_{0,j}|\le C,\quad\|e_j\|_{2,\mathrm{avg}}\to0.
\tag{D.7.2}
\]

Hence the lemma forces

\[
25\,a_j^2/V_j\longrightarrow0.
\tag{D.7.3}
\]

But the already established actual zero-endpoint identities give

\[
a_j^2/V_j\longrightarrow\frac1{6\pi^2}>0.
\tag{D.7.4}
\]

Thus (7.1) excludes the zero endpoint. In particular, if (7.1) is
established with constants depending on a fixed finite bound
\(\chi/A\le B\), this proves O10 for \textbf{every fixed finite \(B\)}
after the existing finite-parameter estimate and exact Yamabe transfer.
It would replace the earlier finite threshold \(39/4+\delta\) by
arbitrary fixed density bounds.

This note does \textbf{not} claim (7.1) without a bound on the negative
Weyl budget. An estimate whose right side grows with \(\fint|W^-|^2\) is
not the uniform bound required by the lemma along an unrestricted
sequence.

\hypertarget{what-the-older-rank-two-argument-could-not-see}{%
\subsection{8. What the older rank-two argument could not
see}\label{what-the-older-rank-two-argument-could-not-see}}

For a rank-two section, Chern localization involves a fourth-power
derivative term. Fixed-parameter smoothing and an amplitude-exceptional
set of measure \(O(\varepsilon)\) do not make its fourth-power energy
vanish. For a line section, the topological pairing in (6.1) involves
only the \textbf{quadratic} derivative expression \(d\chi\wedge\alpha\),
paired with a bounded curvature model. The logarithmic resolvent choice
upgrades its estimate from \(O(1)\) to \(o(1)\).

This distinction is the mathematical reason to pass from the opposite
rank-two torsion bundle to the third-torsion line. It is not an
assertion that rank-two characteristic energy or the existing Weyl tail
disappears automatically.

\end{document}
